\section{Methods and Tools}

This section outlines the methodological approach, tools, and experimental setup used throughout the security assessment.

\subsection{Experimental Environment}

\subsubsection{Hardware and Platform}
All testing was conducted on JMU Computer Science laboratory machines running Ubuntu/Kali Linux, accessed remotely through SSH. The hardware specifications included (see Figure~\ref{fig:A-misc-16-system-info-kali-linux-apple-silicon}):
\begin{itemize}
    \item Architecture: ARM64 (aarch64) processors (Apple Silicon hosts in lab)
    \item Multi-core CPU configuration optimized for parallel processing
    \item Sufficient RAM for wordlist operations and large file processing
    \item High-speed network connectivity for VPN-based web testing
\end{itemize}

\subsubsection{Software Environment}
The testing environment utilized Kali Linux 2023.3 with the following key components (tool versions referenced per-task in Results):
\begin{itemize}
    \item Linux kernel 6.12.38+kali-arm64 with ARM optimizations
    \item Python 3.x for custom analysis scripts
    \item OpenVPN client for secure access to target web servers
    \item Various security analysis tools (detailed below)
\end{itemize}

\subsection{Security Tools and Frameworks}

\subsubsection{Password Hash Cracking Tools}

\paragraph{\ToolName{Ophcrack}}
Ophcrack specializes in LM hash cracking using precomputed rainbow tables. Key characteristics:
\begin{itemize}
    \item Utilizes time-space trade-off for rapid hash lookups
    \item Requires minimal computational resources during attack phase
    \item Highly effective against LM hashes due to their structural weaknesses
    \item Downloaded rainbow tables from \url{https://ophcrack.sourceforge.io/tables.php}
\end{itemize}

\paragraph{\ToolName{John the Ripper}}
A versatile password cracker supporting multiple hash formats:
\begin{itemize}
    \item Version 1.9.0-jumbo-1+bleeding-562f1d5 with ARM optimizations
    \item MD4 128/128 ASIMD 4x2 algorithm implementation for NTLM hashes
    \item MKPC (Most Keys Per Crypt): 4096 for optimal performance
    \item Support for wordlist attacks, rule-based mutations, and incremental mode
\end{itemize}

\paragraph{\ToolName{Hydra}}
Online password guessing tool for web application testing:
\begin{itemize}
    \item Support for multiple protocols (HTTP forms, Basic Auth, Digest Auth)
    \item Configurable rate limiting and session management
    \item Integration with custom wordlists and credential sets
    \item Detailed logging and progress monitoring capabilities
\end{itemize}

\subsubsection{Supporting Tools}
\begin{itemize}
    \item \textbf{OpenSSL}: Certificate and private key analysis
    \item \textbf{Custom Python Scripts}: Attack visualization and statistical analysis
    \item \textbf{Wordlist Management}: pw-inspector for length filtering, custom wordlist generation
\end{itemize}

\subsection{Methodological Approach}

\subsubsection{Strategic Security Analysis}
Following the project's emphasis on developing a security mindset, the assessment prioritized:

\begin{enumerate}
    \item \textbf{Weakest Link Identification}: Analyzing each system to identify the most vulnerable authentication components
    \item \textbf{Resource Optimization}: Implementing time-limited attacks with specific yield thresholds
    \item \textbf{ROI-Focused Testing}: Prioritizing attack vectors based on success probability versus computational cost
    \item \textbf{Comprehensive Documentation}: Maintaining detailed logs of all commands, timing, and results
\end{enumerate}

\subsubsection{Attack Phases}

\paragraph{Phase 1: Reconnaissance and Analysis}
\begin{itemize}
    \item Password Verification Data (PVD) file format analysis
    \item Hash type identification and field structure examination
    \item Tool capability assessment and attack planning
\end{itemize}

\paragraph{Phase 2: Offline Dictionary Attacks}
\begin{itemize}
    \item Wordlist preparation and optimization
    \item Sequential tool deployment based on hash vulnerability
    \item Real-time monitoring and success tracking
\end{itemize}

\paragraph{Phase 3: Online Password Guessing}
\begin{itemize}
    \item VPN connectivity establishment to target networks
    \item Web application authentication mechanism analysis
    \item Controlled brute-force testing with rate limiting
\end{itemize}

\paragraph{Phase 4: Token Authentication Analysis}
\begin{itemize}
    \item RSA certificate and private key extraction
    \item Cryptographic parameter identification
    \item Security implementation assessment
\end{itemize}

\subsection{Data Collection and Analysis}

All experimental activities were comprehensively logged, including:
\begin{itemize}
    \item Complete command histories with timestamps
    \item Tool output files with success/failure indicators
    \item Runtime measurements and performance metrics
    \item Attack success rates and pattern analysis
    \item Security vulnerability assessments and recommendations
\end{itemize}

Evidence collection followed strict protocols to ensure reproducibility and academic integrity, with all results saved to structured directories for subsequent analysis and report generation.